\documentclass[11pt,a4paper]{article}
\usepackage[margin=1in]{geometry}
\usepackage{hyperref}
\usepackage{enumitem}
\usepackage{graphicx}
\usepackage{xcolor}
\usepackage{float}
\usepackage{longtable}
% Tighten spacing around figures so they don't leave large gaps
\setlength{\textfloatsep}{4pt plus 1pt minus 1pt}
\setlength{\floatsep}{4pt plus 1pt minus 1pt}
\setlength{\intextsep}{4pt plus 1pt minus 1pt}
\setlength{\abovecaptionskip}{3pt}
\setlength{\belowcaptionskip}{3pt}
\usepackage{booktabs}
\usepackage{caption}

% -----------------------------------------------------
% bvraghav-tiet-ucs503p-article.sty
% -----------------------------------------------------

% Variable: Lab Instructor
\makeatletter \newcommand{\BvrSetLabInstructor}[1]{%
  \gdef\Bvr@LabInstructor{#1}}
% default
\newcommand{\Bvr@LabInstructor}{Dr. Raghav %
  B. Venkataramaiyer}
% public accessor
\newcommand{\BvrLabInstructorValue}{\Bvr@LabInstructor}
\makeatother

% Add subtitle
\usepackage{titling}
% -- Set title bold
\pretitle{\begin{center}
  
\includegraphics[width=3.5cm]{thapar_logo.png}\\[0.8em]
  \bfseries\fontsize{18bp}{18bp}\selectfont}
\makeatletter
\newcommand{\BvrSubtitle}[1]{%
  \posttitle{%
    \par\end{center}
    \begin{center}\large#1\end{center}
    \vskip 0.5em}%
}
\makeatother

% Hints in the section title(s)
\newcommand{\BvrHint}[1]{%
  {\normalfont\normalsize\itshape\hfill #1}}

% -----------------------------------------------------

\title{Athena: An AI-Powered Campus Query Bot}

\BvrSetLabInstructor{Ms. Jhonsy Bansal}

\BvrSubtitle{%
  Project Proposal for UCS503 \\
  Team: The 3 Stooges \\
  Submitted to: \BvrLabInstructorValue \\ \bfseries
  Thapar Institute of Engineering and Technology%
}

\author{
  Armandeep Singh Taggar, CSED \\
  \small Roll No: \texttt{1024030219}
  \\[0.5em]
  Sanchit Narula, CSED \\
  \small Roll No: \texttt{1024031048}
  \\[0.5em]
  Nipun Gupta, CSED \\
  \small Roll No: \texttt{1024031169}
}

\date{\today}

\begin{document}

\maketitle

\section{Higher-order goal%
  \BvrHint{\ldots secondary framing}}
This project aims to improve the student experience during and after
college admission by reducing the time and effort required to find
accurate campus information. While the underlying problem is broader
than any single institution, the immediate engineering objective is to
translate \emph{centralized, verified access to campus knowledge} into
a measurable, deployable software solution.

\section{Time-to-value %
\BvrHint{\ldots primary heuristic}}
\begin{itemize}[leftmargin=0.7cm]
    \item We will deliver meaningful increments quickly by prototyping the core query-bot workflow (ask a question, retrieve from a verified knowledge base, get a grounded answer) and validating it within a short iteration window.
    \item We will prioritize fast feedback over perfection by using weekly iterations to reduce release friction, starting with a single-college deployment before generalizing to multiple colleges.
    \item Success is defined by what can be delivered and measured in the first iteration (a working query bot on institutional data), then improved based on user feedback.
\end{itemize}

\section{Problem Statement}
Having experienced this ourselves, we identified a common challenge
faced by students both before and after joining college: accessing
accurate, timely, and centralized campus information. Common pain
points include:
\begin{itemize}[leftmargin=0.7cm]
    \item \textbf{Fragmentation:} official notices, events, venues, society information, and academic details are scattered across multiple websites, portals, PDFs, and social media pages.
    \item \textbf{Low discoverability:} students cannot easily search or navigate this scattered information, especially newcomers unfamiliar with the campus.
    \item \textbf{No conversational access:} existing sources require manual searching rather than allowing a student to simply ask a question in natural language.
    \item \textbf{No verification layer:} there is no reliable way to confirm that information (or the person accessing it) genuinely belongs to a given institution.
\end{itemize}

\textbf{Why this matters now (contextual relevance):} As colleges publish more information digitally but across disconnected channels, even a modest centralization and search layer can noticeably reduce the time students spend hunting for basic information.

\section{Proposed Solution %
\BvrHint{\ldots Software Proposition}}
\subsection{Overview}
Build a \textbf{web-based}, multi-college platform, ``Athena'', a
centralized AI-powered campus assistant with the following components:
\begin{itemize}[leftmargin=0.7cm]
    \item \textbf{AI Campus Query Bot:} a natural-language interface for students to ask questions (e.g., ``What is the hostel admission procedure?'') and receive answers grounded in verified institutional data.
    \item \textbf{Navigation:} an interactive 2D campus map (with a roadmap toward 3D) to help students locate classrooms, departments, facilities, and event venues.
    \item \textbf{Academics:} centralized course codes, faculty details, syllabi, timetables, date sheets, and a CGPA calculator.
    \item \textbf{Societies \& Events:} consolidated information on student societies, upcoming events, and announcements.
    \item \textbf{Personalized Dashboard \& Reminders:} a student-specific dashboard for managing reminders and important dates.
    \item \textbf{Admin access \& student verification:} each institution maintains its own verified knowledge base, with a mechanism to confirm a student belongs to the selected institution.
\end{itemize}

\subsection{Core workflow}
\begin{enumerate}[leftmargin=0.7cm]
    \item Student selects/registers with their institution and is verified as belonging to it.
    \item Student asks a question in natural language through the query bot.
    \item The system identifies the student's institution and retrieves relevant context from that institution's verified knowledge base (semantic search over document embeddings).
    \item The retrieved context is passed to the LLM, which generates a clear, contextual response grounded in the retrieved data.
    \item For location-based queries, the response is additionally resolved into structured data (building, floor, coordinates) and visually surfaced on the campus map.
\end{enumerate}

\subsection{Operational constraints for an educational, multi-college setting}
\begin{itemize}[leftmargin=0.7cm]
    \item Keep the solution usable on low-to-mid range devices via responsive web design.
    \item Ground every answer in retrieved, verified institutional data rather than relying on a general-purpose LLM alone, since general LLMs have no access to a college's latest notices, room allocations, or internal rules.
    \item Design each institution's knowledge base and admin access independently, so onboarding a new college does not require changing the underlying platform.
\end{itemize}

\section{Solution Approach %
\BvrHint{\ldots Engineering focus}}
This is an engineering course project, so we focus on system-level
integration of existing, mature components rather than novel algorithm
research.

\subsection{Retrieval-grounded answering %
\BvrHint{\ldots non-research}}
Answer generation is retrieval-augmented rather than a research
contribution in itself:
\begin{itemize}[leftmargin=0.7cm]
    \item Documents (notices, syllabi, policies, etc.) are converted into vector embeddings using Sentence Transformers and stored via pgvector.
    \item A student query is embedded and matched against this store using semantic similarity, without claiming state-of-the-art retrieval.
    \item The retrieved context, not the raw query alone, is passed to the LLM (GPT-4o-mini) so answers stay grounded and easy to update --- when college data changes, only the knowledge base needs updating, not the model.
\end{itemize}

\subsection{Web architecture %
\BvrHint{\ldots web-based solution}}
A typical multi-tier web architecture:
\begin{itemize}[leftmargin=0.7cm]
    \item \textbf{Frontend:} React.js (Vite) with Tailwind CSS for the student/admin interface, dashboards, and 2D map (with Three.js planned for interactive 3D navigation).
    \item \textbf{Backend API:} Django + DRF for authentication, admin management, and the API layer connecting the frontend to the RAG pipeline.
    \item \textbf{Database:} PostgreSQL with the pgvector extension to store students, colleges, courses, faculty, events, notices, and embeddings.
    \item \textbf{Document processing:} PyMuPDF to extract text from PDFs and college documents feeding the knowledge base.
\end{itemize}

\subsection{CI/CD and fast delivery enabler}
CI/CD will be used to:
\begin{itemize}[leftmargin=0.7cm]
    \item Automatically build and run tests on every change (Git + GitHub).
    \item Produce repeatable deployment artifacts to staging (Render/Railway for backend, Vercel for frontend).
    \item Enable safe and frequent releases of incremental features as the knowledge base and feature set grow.
\end{itemize}

\section{Evaluation Criterion %
\BvrHint{\ldots measurable and attributable}}
We define success using measurable metrics tied to the core query-bot
and navigation workflow.

\subsection{Primary evaluation metric}
\textbf{Answer groundedness \& relevance:} the proportion of query-bot
responses that are factually grounded in the retrieved institutional
context and judged relevant by the asking student.
\begin{itemize}[leftmargin=0.7cm]
    \item Target: high grounded-answer rate on a curated RAG evaluation test dataset during the pilot window.
    \item Attribution: measured via RAG evaluation methodology comparing generated answers against the source documents retrieved.
\end{itemize}

\subsection{Secondary evaluation metrics}
\begin{itemize}[leftmargin=0.7cm]
    \item \textbf{Query response time:} end-to-end latency from question submission to answer display.
    \item \textbf{Navigation accuracy:} whether the correct building/floor/room is displayed/highlighted for location queries.
    \item \textbf{Student verification success rate:} percentage of legitimate students correctly verified against their institution.
    \item \textbf{Reliability:} availability of staging/production for login and querying during business hours.
\end{itemize}

\subsection{Pilot validation plan %
\BvrHint{\ldots high-level}}
\begin{itemize}[leftmargin=0.7cm]
    \item Recruit a pilot group of students (and an admin/staff role for knowledge-base maintenance) from one institution to begin with.
    \item Compare metrics before/after within the iteration window by logging query outcomes and gathering direct student feedback on answer usefulness.
\end{itemize}

\section{Scalability %
\BvrHint{\ldots with foresight}}
The proposition should scale in both theoretical and practical senses.

\begin{enumerate}[leftmargin=0.7cm]
    \item \textbf{Scalable (theoretically):} The system should support growth in number of institutions, documents, and concurrent users by:
    \begin{itemize}[leftmargin=0.7cm]
        \item decoupling components (frontend/backend/RAG pipeline),
        \item indexing embeddings via pgvector for efficient semantic search at scale,
        \item using a stateless API design so institutions can be onboarded without platform changes.
    \end{itemize}

    \item \textbf{Operationally deployable (practically):} The system should run on common web hosting:
    \begin{itemize}[leftmargin=0.7cm]
        \item managed PostgreSQL database with pgvector,
        \item platform-hosted deployment (Render/Railway, Vercel),
        \item CI/CD pipeline for staging/production.
    \end{itemize}
\end{enumerate}

\section{Engine availability heuristic}
A mature set of ``engines'' (tooling/services) is available to
implement this as a web-based project:
\begin{itemize}[leftmargin=0.7cm]
    \item Web framework for REST APIs (Django + DRF) and a reactive frontend (React.js/Vite).
    \item Managed relational database with vector support (PostgreSQL + pgvector).
    \item Embedding model (Sentence Transformers) and LLM (OpenAI GPT-4o-mini) for the RAG pipeline.
    \item Document-processing library (PyMuPDF) for ingesting PDFs into the knowledge base.
    \item Authentication approach (Django Auth) and institutional-email-based student verification.
    \item Test framework (PyTest) and API testing tool (Postman) for unit/integration tests.
    \item CI runner and staging deployment targets (Render/Railway, Vercel).
\end{itemize}
We will use these mature components to focus on workflow design,
grounded-answer correctness, and measurement rather than inventing new
infrastructure.

\subsection{Tech Stack / Tools Required}
\textbf{v1: tech we are using currently}

\vspace{0.3cm}

\renewcommand{\arraystretch}{1.2}

\begin{longtable}{|p{2.5cm}|p{3.5cm}|p{6cm}|p{1.5cm}|}

\hline
\textbf{Layer} &
\textbf{Technology} &
\textbf{Purpose} &
\textbf{Status} \\
\hline

Frontend &
React.js (Vite) &
Build the student/admin interface and dashboards &
v1 \\
\hline

Styling &
Tailwind CSS &
UI structure and styling &
v1 \\
\hline

Campus Navigation &
2D map (SVG, floor-plan images) &
Locate classrooms, departments, event venues &
v1 \\
\hline

Campus Navigation &
Three.js &
Interactive 3D campus navigation &
roadmap \\
\hline

Backend &
Django + DRF &
Core logic, auth, admin management, API layer &
v1 \\
\hline

Database &
PostgreSQL &
Store students, colleges, courses, faculty, events, notices &
v1 \\
\hline

Vector search &
pgvector &
Store embeddings, semantic similarity search for RAG &
v1 \\
\hline

Embeddings &
Sentence Transformers &
Convert documents and queries into vector embeddings &
v1 \\
\hline

RAG pipeline &
Custom Python pipeline &
Retrieve documents, construct context for the LLM &
v1 \\
\hline

LLM &
OpenAI GPT-4o-mini &
Generate natural-language answers from retrieved context &
v1 \\
\hline

Document processing &
PyMuPDF &
Extract text from PDFs and college documents &
v1 \\
\hline

Authentication &
Django Auth &
Secure student and admin authentication &
v1 \\
\hline

Student verification &
Institutional email check &
Verify student belongs to selected institution &
v1 \\
\hline

Admin management &
Django Admin &
Maintain and update institutional information &
v1 \\
\hline

Location data &
Structured JSON &
Building, floor, room mapping for navigation &
v1 \\
\hline

API testing &
Postman &
Test backend APIs during development &
v1 \\
\hline

Version control &
Git + GitHub &
Code management and team collaboration &
v1 \\
\hline

Deployment &
Render/Railway, Vercel &
Host and serve backend, database, and frontend &
v1 \\
\hline

\end{longtable}

\textbf{Roadmap: for future development} --- see Section~\ref{sec:scope} for the future-scope items planned beyond v1.

\section{System Design}
The figures below summarize the intended system design.

\begin{center}
\includegraphics[width=0.65\textwidth]{case_diagram.png}
\captionof{figure}{Case Diagram --- students interact with the query bot, navigation, academics, and reminder features, while admins manage the institutional knowledge base.}
\end{center}

\begin{center}
\includegraphics[width=0.65\textwidth]{architecture.png}
\captionof{figure}{System Architecture --- a student request travels from the React frontend down to the LLM, with the RAG pipeline in the middle grounding answers in verified data.}
\end{center}

\begin{center}
\includegraphics[width=0.7\textwidth]{er_diagram.png}
\captionof{figure}{ER Diagram --- core entities: Institution, Student, Course, Faculty, Event, Society, Reminder, and Knowledge\_Document.}
\end{center}

\begin{center}
\includegraphics[width=0.7\textwidth]{sequence_diagram.png}
\captionof{figure}{Sequence Diagram --- message exchange for a single query across student, frontend, backend, RAG pipeline/vector store, and LLM.}
\end{center}

\begin{center}
\includegraphics[width=0.6\textwidth]{class_diagram.png}
\captionof{figure}{Class Diagram --- core classes and their relationships, mirroring the ER diagram with added behavior relevant to the query bot and reminders.}
\end{center}

\begin{center}
\includegraphics[width=0.65\textwidth]{deployment_diagram.png}
\captionof{figure}{Deployment Diagram --- browser client, Vercel frontend, Render/Railway backend and database, and the external OpenAI API.}
\end{center}

\begin{center}
\includegraphics[width=0.65\textwidth]{workflow.png}
\captionof{figure}{Athena Work Flow.}
\end{center}

\section{Project scope and deliverables %
\BvrHint{\ldots iteration-friendly}}
\label{sec:scope}
\subsection{Initial deliverable %
\BvrHint{\ldots first iteration}}
\begin{itemize}[leftmargin=0.7cm]
    \item Student/admin frontend (React + Tailwind) and Django + DRF backend with authentication
    \item Institutional knowledge base ingestion pipeline (PyMuPDF $\rightarrow$ embeddings $\rightarrow$ pgvector)
    \item AI Campus Query Bot answering questions grounded in retrieved context (RAG with GPT-4o-mini)
    \item 2D campus map for locating classrooms, departments, and event venues
    \item Student verification via institutional email check
    \item CI: build + backend unit tests; CD to staging (Render/Railway, Vercel)
\end{itemize}

\subsection{Subsequent deliverables}
\begin{itemize}[leftmargin=0.7cm]
    \item Interactive 3D campus navigation (Three.js + Blender models) with turn-by-turn direction
    \item Personalized dashboard and reminders for academic/campus dates
    \item Admin dashboard for knowledge-base maintenance across multiple colleges
    \item Voice assistant, mobile app, university ERP and attendance integration, personalized academic assistant
\end{itemize}

\section{Risks and mitigations}
\begin{itemize}[leftmargin=0.7cm]
    \item \textbf{Ungrounded/incorrect answers:} rely on RAG evaluation with a test dataset before release; never let the LLM answer without retrieved context for institution-specific questions.
    \item \textbf{Student verification bypass:} keep verification (institutional email check) mandatory before institution-specific data is served.
    \item \textbf{Knowledge-base staleness:} give admins a straightforward Django Admin workflow to update documents so the vector store reflects current information.
    \item \textbf{Scaling to multiple colleges:} keep institution data and knowledge bases logically separated from the start so onboarding a new college does not require re-architecture.
\end{itemize}

\section{Summary}
This project proposes Athena, a deployable, multi-college AI campus
assistant that centralizes fragmented campus information behind a
single natural-language interface. It meets the course criteria by
emphasizing time-to-value through rapid iterations starting with a
single-college pilot, implementing CI/CD to support frequent safe
releases, and using measurable evaluation criteria centered on
grounded, relevant answers. Its contribution lies in system-level
integration --- verified institutional knowledge, student
verification, semantic search, retrieval-augmented generation, and an
interactive campus map --- rather than research-driven novelty in any
single component.

\end{document}